% =====================================================================
%  ARRIVAL-MNEMO: Persistent Memory Architecture
%  Companion Paper to the ARRIVAL Protocol
%  DOI: 10.5281/zenodo.18893515
%  License: CC BY-NC 4.0 (text), AGPL-3.0-or-later (code)
%  Author: Mefodiy Kelevra — ORCID 0009-0003-4153-392X
% =====================================================================
\documentclass[11pt,a4paper]{article}

% ── Geometry ──────────────────────────────────────────────────────────
\usepackage[margin=2.4cm,top=2.8cm,bottom=2.8cm]{geometry}

% ── Fonts & Encoding ─────────────────────────────────────────────────
\usepackage[T1]{fontenc}
\usepackage[utf8]{inputenc}
\usepackage{lmodern}
\usepackage{microtype}

% ── Math ─────────────────────────────────────────────────────────────
\usepackage{amsmath,amssymb,amsthm}

% ── Tables & Figures ─────────────────────────────────────────────────
\usepackage{booktabs}
\usepackage{tabularx}
\usepackage{multirow}
\usepackage{graphicx}
\usepackage{float}

% ── Colours & Boxes ──────────────────────────────────────────────────
\usepackage[dvipsnames]{xcolor}
\usepackage[most]{tcolorbox}

% ── Code listings ────────────────────────────────────────────────────
\usepackage{listings}
\lstset{
  basicstyle=\small\ttfamily,
  keywordstyle=\color{NavyBlue}\bfseries,
  commentstyle=\color{OliveGreen},
  stringstyle=\color{BrickRed},
  breaklines=true,
  frame=single,
  framerule=0.4pt,
  xleftmargin=6pt,
  xrightmargin=6pt,
  language=Python,
  showstringspaces=false,
}

% ── Hyperlinks ───────────────────────────────────────────────────────
\usepackage[colorlinks=true,linkcolor=NavyBlue,citecolor=OliveGreen,
            urlcolor=BrickRed]{hyperref}
\usepackage{url}

% ── Bibliography ─────────────────────────────────────────────────────
\usepackage[numbers,sort&compress]{natbib}

% ── Section formatting ───────────────────────────────────────────────
\usepackage{titlesec}
\titleformat{\section}{\Large\bfseries}{\thesection.}{0.5em}{}
\titleformat{\subsection}{\large\bfseries}{\thesubsection}{0.5em}{}
\titleformat{\subsubsection}{\normalsize\bfseries}{\thesubsubsection}{0.5em}{}

% ── Theorem-like environments (tcolorbox style) ─────────────────────
\newtcolorbox{theorembox}[1][]{
  colback=blue!3, colframe=NavyBlue!70!black,
  fonttitle=\bfseries, title={#1},
  boxrule=0.6pt, arc=2pt, left=6pt, right=6pt, top=4pt, bottom=4pt
}
\newtcolorbox{resultbox}[1][]{
  colback=red!3, colframe=BrickRed!60!black,
  fonttitle=\bfseries, title={#1},
  boxrule=0.6pt, arc=2pt, left=6pt, right=6pt, top=4pt, bottom=4pt
}
\newtcolorbox{definitionbox}[1][]{
  colback=green!3, colframe=OliveGreen!70!black,
  fonttitle=\bfseries, title={#1},
  boxrule=0.6pt, arc=2pt, left=6pt, right=6pt, top=4pt, bottom=4pt
}
\newtcolorbox{alertbox}[1][]{
  colback=yellow!5, colframe=Dandelion!80!black,
  fonttitle=\bfseries, title={#1},
  boxrule=0.6pt, arc=2pt, left=6pt, right=6pt, top=4pt, bottom=4pt
}

% ── Misc ─────────────────────────────────────────────────────────────
\usepackage{enumitem}
\setlist{nosep,leftmargin=*}
\usepackage{xspace}
\newcommand{\arrival}{\textsc{Arrival}\xspace}
\newcommand{\mnemo}{\textsc{Arrival-Mnemo}\xspace}
\newcommand{\care}{\textsc{Care}\xspace}
\newcommand{\pp}{\,pp\xspace}

% ── Header / Footer ──────────────────────────────────────────────────
\usepackage{fancyhdr}
\pagestyle{fancy}
\fancyhf{}
\fancyhead[L]{\small\textit{ARRIVAL-MNEMO — Memory Architecture}}
\fancyhead[R]{\small\textit{DOI: 10.5281/zenodo.18893515}}
\fancyfoot[C]{\thepage}
\renewcommand{\headrulewidth}{0.4pt}

% =====================================================================
\begin{document}

% ── Title Page ───────────────────────────────────────────────────────
\begin{titlepage}
\centering
\vspace*{2cm}

{\Huge\bfseries ARRIVAL-MNEMO}\\[0.6cm]
{\Large Persistent Memory Architecture for\\
Cross-Architecture AI Coordination}\\[2.5cm]

{\large Mefodiy Kelevra}\\[0.3cm]
{\small ORCID:
  \href{https://orcid.org/0009-0003-4153-392X}{0009-0003-4153-392X}}\\[2cm]

{\normalsize March 4, 2026 \enspace|\enspace Version 3.0}\\[0.5cm]
{\small DOI: \href{https://doi.org/10.5281/zenodo.18893515}%
  {10.5281/zenodo.18893515}}\\[0.3cm]
{\small Mathematical Foundation: MEANING-CRDT v1.1
  (DOI: \href{https://doi.org/10.5281/zenodo.18702383}%
  {10.5281/zenodo.18702383})}\\[2cm]

\rule{0.6\textwidth}{0.4pt}\\[0.8cm]

{\footnotesize
  \textbf{License} — Text: CC BY-NC 4.0 \enspace|\enspace
  Code: AGPL-3.0-or-later\\[0.4cm]
  Companion to: \arrival Protocol v2.0
  (DOI: \href{https://doi.org/10.5281/zenodo.18893515}%
  {10.5281/zenodo.18893515})\\[0.2cm]
  Repository:
  \url{https://github.com/mifkilla/Arrival-Protocol}
}

\vfill
\end{titlepage}

% ── Table of Contents ────────────────────────────────────────────────
\tableofcontents
\newpage

% =====================================================================
\section{Abstract}
% =====================================================================

The \arrival Protocol enables cross-architecture AI-to-AI coordination
through structured semantic @-atoms (DEUS.PROTOCOL~v0.5), validated
across 2,200+ experiments.  However, a critical limitation persists: all
coordination knowledge is lost between sessions.  We present \mnemo, a
persistent memory architecture extending the protocol with four memory
layers (Episodic, Procedural, Semantic, Meta), a \texttt{@MEMORY.*} atom
family, and MEMORY-CRDT~v1.1 --- a formal specification for
cross-temporal memory synchronization using the same \care Resolve
mathematics that powers real-time dialogue convergence.

We conducted three experimental phases.  Phase~14 reveals a
\textbf{negative result}: naive keyword-based global memory injection
DEGRADES accuracy by $-5.7$\pp (McNemar $p = 0.724$) on GPQA Diamond,
discovering the \textit{hypercorrection effect} in LLM ensembles --- a
cognitive bias previously documented only in humans.  Phase~15 introduces
\care-ALERT, a gated real-time intervention system that restores
baseline accuracy while achieving the first statistically significant
improvement in all \arrival experiments: \care Resolve is significantly
higher than under global injection (Mann--Whitney $p = 0.042$).

Key findings: (1)~memory creation via auto-logging is solved;
(2)~global pre-session injection triggers hypercorrection;
(3)~gated real-time intervention avoids hypercorrection while
improving consensus quality; (4)~specific memory helps but generic
memory hurts AI coordination.

\textbf{Keywords:} AI coordination, persistent memory, CRDT, semantic
atoms, multi-agent systems, hypercorrection, CARE-ALERT

% =====================================================================
\section{Introduction}\label{sec:intro}
% =====================================================================

\subsection{The Amnesia Problem}

Modern AI coordination protocols operate in a fundamentally amnesiac
mode.  When a multi-agent session ends, all accumulated knowledge
vanishes.  The agents retain no recollection of which strategies worked,
which information sources proved reliable, or how their peers typically
behave.  Every new session starts from absolute zero.

The \arrival Protocol~\cite{kelevra2026arrival} demonstrates that
cross-architecture AI agents can achieve structured coordination through
semantic @-atoms, with results validated across 2,200+ experiments,
23~phases, and 17 LLM architectures.  But this coordination knowledge
exists only in log files, inaccessible to future sessions.

\subsection{Research Questions}

\textbf{RQ1}: What memory architecture enables persistent learning
across sessions while maintaining compatibility with the @-atom
communication framework?

\textbf{RQ2}: Can the CRDT mathematics used for real-time dialogue
convergence (\care Resolve) be extended to cross-temporal memory
synchronization?

\textbf{RQ3}: Does memory injection improve coordination quality
compared to memoryless baselines?

\subsection{Contributions}

\begin{enumerate}
  \item \textbf{Four-layer memory taxonomy} (Episodic, Procedural,
    Semantic, Meta) with \texttt{@MEMORY.*} atoms extending
    DEUS.PROTOCOL~v0.5.
  \item \textbf{MEMORY-CRDT v1.1}: Formal specification for memory
    synchronization using utility-weighted LWW-Register CRDT with
    convergence guarantees.
  \item \textbf{Hypercorrection discovery}: Global pre-session memory
    injection degrades LLM ensemble performance --- a cognitive bias
    previously documented only in humans.
  \item \textbf{CARE-ALERT system}: Gated real-time intervention
    achieving the first statistically significant improvement across
    all \arrival experiments (Mann--Whitney $p = 0.042$).
\end{enumerate}

% =====================================================================
\section{Related Work}\label{sec:related}
% =====================================================================

\textbf{Retrieval-Augmented Generation} (Lewis et al.,
2020)~\cite{lewis2020rag} retrieves relevant documents to augment
generation but operates at the document level without structured memory
layers or cross-agent synchronization.

\textbf{MemGPT} (Packer et al., 2023)~\cite{packer2023memgpt}
introduces hierarchical memory management for LLMs with main context
and external storage but targets single-agent scenarios.

\textbf{Generative Agents} (Park et al., 2023)~\cite{park2023agents}
implement memory streams for simulated social agents with reflection
and planning but focus on narrative coherence rather than structured
coordination protocols.

\textbf{MEANING-CRDT v1.1}~\cite{kelevra2026crdt} introduced \care
Resolve for real-time dialogue synchronization.  MEMORY-CRDT extends
this mathematical framework to the temporal dimension.

\textbf{CRDTs} (Shapiro et al., 2011)~\cite{shapiro2011crdt} provide
formal guarantees for distributed data convergence.  Our adaptation
replaces physical timestamps with utility scores.

% =====================================================================
\section{Architecture}\label{sec:arch}
% =====================================================================

\subsection{Four Memory Layers}

\mnemo organizes persistent knowledge into four layers:

\begin{definitionbox}[Memory Layers]
\begin{description}[style=nextline]
  \item[Episodic] Session-level records: task, models, outcome, \care
    score. TTL: 7--30 days. \textit{``What happened.''}
  \item[Procedural] Validated strategies with success rates.
    Persistence requires $n_{\text{trials}} \geq 3$.
    \textit{``How to act.''}
  \item[Semantic] Derived knowledge: rules, facts, calibration data.
    Persistence requires $\text{evidence\_count} \geq 3$.
    \textit{``What is true.''}
  \item[Meta] Agent calibration: trust scores, domain-specific accuracy,
    behavioral patterns. Updated via exponential moving average.
    \textit{``Who to trust.''}
\end{description}
\end{definitionbox}

\subsection{@MEMORY.* Atom Family}

Nine new atoms extend DEUS.PROTOCOL~v0.5:

\begin{table}[H]
\centering\small
\caption{\texttt{@MEMORY.*} atom family.}\label{tab:memoryatoms}
\begin{tabular}{@{}ll@{}}
\toprule
\textbf{Atom} & \textbf{Function} \\
\midrule
\texttt{@MEMORY.EPISODIC}   & Reference/inject session-level memory \\
\texttt{@MEMORY.PROCEDURAL} & Reference/inject strategy-level memory \\
\texttt{@MEMORY.SEMANTIC}   & Reference/inject knowledge-level memory \\
\texttt{@MEMORY.META}       & Reference/inject calibration data \\
\texttt{@MEMORY.INJECT}     & Request memory injection during session \\
\texttt{@MEMORY.STORE}      & Request memory persistence after session \\
\texttt{@MEMORY.FORGET}     & Flag memory for eviction \\
\texttt{@MEMORY.CONFLICT}   & Signal contradicting memories \\
\texttt{@MEMORY.MERGE}      & Request memory consolidation \\
\bottomrule
\end{tabular}
\end{table}

\subsection{Injection Protocol}

Memory injection follows a seven-step pipeline:
\begin{enumerate}
  \item \textbf{LOAD}: Read memory store from persistent storage (JSON).
  \item \textbf{FILTER}: Select memories relevant to current
    \texttt{@GOAL} using keyword matching (MVP) or embedding similarity.
  \item \textbf{FORMAT}: Convert to \texttt{@MEMORY.*} atom notation.
  \item \textbf{INJECT}: Append to agent system prompts as
    \texttt{[MEMORY CONTEXT]} block.
  \item \textbf{RUN}: Execute \arrival dialogue with memory-augmented
    agents.
  \item \textbf{EXTRACT}: Create new Episodic memories; update Meta
    with calibration data.
  \item \textbf{MERGE}: Integrate via MEMORY-CRDT.
\end{enumerate}

% =====================================================================
\section{MEMORY-CRDT v1.1}\label{sec:crdt}
% =====================================================================

\subsection{From Dialogue to Memory}

MEANING-CRDT merges divergent agent positions within a dialogue:
\[
  \hat{\mathbf{v}} = \frac{\sum_{i} w_i \cdot \mathbf{v}_i}
                          {\sum_{i} w_i}
\]

MEMORY-CRDT applies the same principle across time.  When two memory
stores conflict, the merge operation selects the higher-utility version:
\[
  \hat{m} = \arg\max_{m \in \{m_1, m_2\}} \mu(m)
\]
where $\mu(m)$ is the utility function:
\[
  \mu(m) = 0.25 \cdot \text{recency} + 0.30 \cdot \text{frequency}
         + 0.25 \cdot \text{quality} + 0.20 \cdot \text{validation}
\]

\subsection{Utility Components}

\begin{itemize}
  \item $\text{recency}(m) = \max(0,\; 1.0 - \text{age\_days}/90)$.
    Linear decay over 90-day window.
  \item $\text{frequency}(m) = \min(1.0,\;
    \text{access\_count}/N_{\text{layer}})$.
    $N$ varies by layer: Episodic${}=10$, Procedural${}=20$,
    Semantic${}=5$, Meta${}=10$.
  \item $\text{quality}(m)$: layer-specific (care\_resolve, success\_rate,
    confidence, trust\_score).
  \item $\text{validation}(m)$: evidence strength, layer-dependent.
\end{itemize}

\subsection{CRDT Properties}

\begin{theorembox}[MEMORY-CRDT Properties]
The merge operation is a \textbf{Last-Writer-Wins Register}
(LWW-Register) CRDT where the ``timestamp'' is replaced by utility
score:
\begin{itemize}
  \item \textbf{Commutativity}: $\text{MERGE}(S_1, S_2) =
    \text{MERGE}(S_2, S_1)$.
  \item \textbf{Associativity}: $\text{MERGE}(\text{MERGE}(S_1, S_2),
    S_3) = \text{MERGE}(S_1, \text{MERGE}(S_2, S_3))$.
  \item \textbf{Idempotency}: $\text{MERGE}(S, S) = S$.
\end{itemize}
\end{theorembox}

\subsection{Memory Debt}

Analogous to Meaning Debt:
\[
  \text{MmD} = \frac{\sum_{i} \mu_i \cdot \text{conflict}(m_i)}
                     {\sum_{i} \mu_i}
\]
where $\text{conflict}(m) \in \{0.0, 0.5, 1.0\}$.

\subsection{Convergence}

\begin{theorembox}[Memory Convergence Theorem]
Given a sequence of sessions $\{s_1, \ldots, s_n\}$ each producing
memory updates, the memory store converges to a stable state as
$n \to \infty$, provided: (1)~the forgetting function removes memories
with $\mu < \theta$; (2)~merge resolves conflicts via utility
maximization; (3)~consolidation extracts patterns into semantic
memories.
\end{theorembox}

\subsection{Relationship to MEANING-CRDT}

\begin{table}[H]
\centering\small
\caption{MEANING-CRDT vs.\ MEMORY-CRDT.}\label{tab:crdtcomp}
\begin{tabular}{@{}lll@{}}
\toprule
\textbf{Property} & \textbf{MEANING-CRDT} & \textbf{MEMORY-CRDT} \\
\midrule
Domain          & Dialogue positions & Memory entries \\
Temporal scope  & Within a session   & Across sessions \\
Weight function & Coherence          & Utility \\
Merge granularity & Per-round        & Per-session \\
Convergence target & Consensus       & Consistent memory \\
Debt metric     & Meaning Debt       & Memory Debt \\
CRDT type       & GCounter-like      & LWW-Register \\
\bottomrule
\end{tabular}
\end{table}

\textbf{Unifying principle}: Both CRDTs implement weighted convergence
under conflict, applied at different temporal scales.

% =====================================================================
\section{Pilot Experiment (Phase 14a)}\label{sec:pilot}
% =====================================================================

\subsection{Design}

A controlled experiment with 10 MCQ questions across 5 domains.
Models: DeepSeek V3, Llama 3.3 70B, Qwen 2.5 72B.  Three conditions:
Seed, Control (no memory), Treatment (memory injection).

\subsection{Results}

\begin{table}[H]
\centering\small
\caption{Pilot experiment results.}\label{tab:pilot}
\begin{tabular}{@{}lrrr@{}}
\toprule
\textbf{Condition} & \textbf{Accuracy} & \textbf{Avg \care} &
  \textbf{Cost} \\
\midrule
Seed                  & 3/5 (60\%)       & 0.850 & \$0.008 \\
Control (no memory)   & 4/5 (80\%)       & 0.750 & \$0.008 \\
Treatment (memory)    & \textbf{5/5 (100\%)} & \textbf{0.850} & \$0.010 \\
\bottomrule
\end{tabular}
\end{table}

Memory injection corrected the one control error, improved \care by
$+13.3\%$, and increased atom density by $+8.6\%$.  Total: \$0.026.

% =====================================================================
\section{Phase 14: Large-Scale Validation (Negative Result)}
\label{sec:phase14}
% =====================================================================

\subsection{Design}

Scaled to GPQA Diamond --- 40 graduate-level science
questions~\cite{rein2024gpqa}.  Alpha trio: GPT-4o + DeepSeek-V3 +
Llama-3.3-70B-Instruct.  Memory store: 5 episodic, 2 procedural,
3 meta.  Keyword-based retrieval, top-8 injection, 800-char limit.

\textbf{Data leakage guard}: No correct answers, answer letters,
question text, or choices stored.

\subsection{Results}

\begin{resultbox}[Phase 14 — Negative Result]
\begin{tabular}{@{}lrrrl@{}}
\toprule
\textbf{Metric} & \textbf{Phase 13} & \textbf{Phase 14} &
  \textbf{$\Delta$} & \textbf{$p$-value} \\
\midrule
Accuracy (N=35) & 57.1\% & 51.4\% & $-5.7$\pp &
  McNemar $p = 0.724$ \\
\care Resolve & 0.966 & 0.916 & $-0.050$ & M-W $p = 0.985$ \\
Meaning Debt  & 0.196 & 0.433 & $+0.237$ & M-W $p = 0.994$ \\
\bottomrule
\end{tabular}
\end{resultbox}

Memory injection did not improve accuracy.  The trend is negative but
not statistically significant.

\subsection{The Hypercorrection Effect}

Analysis of auto-stored cognitive scars reveals the failure mechanism.
When told ``in past chemistry sessions, DeepSeekV3 was suppressed,''
GPT-4o begins second-guessing its own correct answers on ALL chemistry
questions, not just the specific ones where suppression occurred.  This
parallels the well-documented human cognitive bias: pre-task warnings
about past mistakes increase anxiety and worsen performance.

\textbf{LLM ensembles empirically exhibit this classic human pattern}:
they need correction DURING the process, not warnings BEFORE it.

\subsection{Implications}

\begin{enumerate}
  \item \textbf{Memory creation is solved}: Auto-logging perfectly
    captures who was suppressed, in which domain, and by how much.
  \item \textbf{Memory delivery via global injection fails}: Static
    pre-session injection of generic memories degrades performance.
  \item \textbf{Specific memory helps, generic memory hurts}: Seen
    questions improved ($+20$\pp); unseen degraded because keyword
    retrieval injected noise.
\end{enumerate}

% =====================================================================
\section{Phase 15: Gated CARE-ALERT Validation}\label{sec:phase15}
% =====================================================================

\subsection{Architecture}

Phase~14's negative result motivates a paradigm shift: from global
injection to \textbf{real-time intervention}.

\begin{alertbox}[CARE-ALERT: Monitor $\to$ Detect $\to$ Intervene]
\begin{enumerate}
  \item R1 (Alpha: GPT-4o) $\to$ clean prompt, no memory
  \item R2 (Beta: DeepSeek-V3) $\to$ clean prompt, no memory
  \item \textit{Checkpoint \#1}: $\text{MD} > 0.5$ $\Rightarrow$
    Level~1 \texttt{@CARE.ALERT}
  \item R3 (Gamma: Llama-3.3) $\to$ prompt + alert (if fired)
  \item \textit{Checkpoint \#2}: $\text{MD} > 0.8$ $\Rightarrow$
    Level~2 ESCALATED
  \item R4 (Alpha: GPT-4o) $\to$ prompt + alert (if fired)
\end{enumerate}
\end{alertbox}

\textbf{Alert design}: Operational, not moralistic (``Beta selected~C,
address with \texttt{@EVIDENCE}'' --- NOT ``you failed before'').
Current state only.  Never contains correct answers, historical data,
or trust scores.  Clean system prompt, identical to Phase~13 baseline.

\subsection{Results: Unseen Questions (N=35)}

\begin{table}[H]
\centering\small
\caption{Three-phase comparison on unseen questions.}\label{tab:p15unseen}
\begin{tabular}{@{}lrrrr@{}}
\toprule
\textbf{Metric} & \textbf{P13} & \textbf{P14} & \textbf{P15} &
  \textbf{$\Delta$ P13$\to$P15} \\
\midrule
Accuracy     & 57.1\% & 51.4\% & \textbf{57.1\%} & $+0.0$\pp \\
\care        & 0.966  & 0.916  & \textbf{0.948}  & $-0.018$ \\
Meaning Debt & 0.196  & 0.433  & 0.452            & $+0.256$ \\
\bottomrule
\end{tabular}
\end{table}

Phase~15 restores accuracy to the Phase~13 baseline.  The $-5.7$\pp
degradation from global injection is fully eliminated.

\subsection{Results: Seen Questions (N=5)}

\begin{table}[H]
\centering\small
\caption{Three-phase comparison on seen questions.}\label{tab:p15seen}
\begin{tabular}{@{}lrrr@{}}
\toprule
\textbf{Metric} & \textbf{P13} & \textbf{P14} & \textbf{P15} \\
\midrule
Accuracy     & 20\%  & 40\%  & \textbf{60\%} \\
\care        & 0.539 & 0.635 & \textbf{0.746} \\
Meaning Debt & 2.455 & 1.944 & \textbf{1.602} \\
\bottomrule
\end{tabular}
\end{table}

Monotonic improvement across all three metrics through all three phases.

\subsection{Results: All 40 Questions}

\begin{table}[H]
\centering\small
\caption{Three-phase comparison on all 40 questions.}\label{tab:p15all}
\begin{tabular}{@{}lrrr@{}}
\toprule
\textbf{Metric} & \textbf{P13} & \textbf{P14} & \textbf{P15} \\
\midrule
Accuracy & 52.5\% & 50.0\% & \textbf{57.5\%} \\
\care    & 0.913  & 0.881  & \textbf{0.923} \\
\bottomrule
\end{tabular}
\end{table}

\begin{resultbox}[First Significant Result]
\care Resolve Phase~14 vs.\ Phase~15 (N=40):\\
\textbf{Mann--Whitney $p = 0.042$} (significant at $\alpha = 0.05$).\\
This is the \textbf{first statistically significant result} across all
\arrival experiments.
\end{resultbox}

\subsection{Alert Statistics}

\begin{table}[H]
\centering\small
\caption{CARE-ALERT firing statistics.}\label{tab:alerts}
\begin{tabular}{@{}lr@{}}
\toprule
\textbf{Statistic} & \textbf{Value} \\
\midrule
Total questions            & 40 \\
Questions with alert       & 7 (17.5\%) \\
Accuracy with alert        & 57.1\% (4/7) \\
Accuracy without alert     & 57.6\% (19/33) \\
\care with alert           & 0.968 \\
\care without alert        & 0.914 \\
\bottomrule
\end{tabular}
\end{table}

The alert does not harm accuracy.  \care is HIGHER when alerts fire
(0.968 vs.\ 0.914).

\subsection{Minority Voice Protection}

Three examples of successful minority voice protection:
\begin{itemize}
  \item \textbf{GPQA\_009}: Beta (DeepSeek) correct, Alpha (GPT-4o)
    wrong.  Alert fired $\to$ final answer followed Beta $\to$ CORRECT.
  \item \textbf{GPQA\_038}: Gamma (Llama) correct, Alpha+Beta wrong.
    Alert fired $\to$ final answer followed Gamma $\to$ CORRECT.
  \item \textbf{GPQA\_003}: Beta correct against Alpha.  Alert $\to$
    CORRECT.
\end{itemize}

\subsection{Analysis}

\textbf{Finding 1}: Gated alerts avoid hypercorrection.  Phase~14
injected memory into 100\% of questions $\to$ $-5.7$\pp.  Phase~15
alerted on 17.5\% $\to$ $+0.0$\pp.

\textbf{Finding 2}: First significant \care improvement
(Mann--Whitney $p = 0.042$).

\textbf{Finding 3}: Precision targeting works.  17.5\% alert rate
demonstrates intervention only when needed.

\textbf{Finding 4}: Memory delivery was the problem, not memory itself:
\begin{itemize}
  \item Phase~13: Baseline (stateless) $\to$ 57.1\%
  \item Phase~14: Global injection $\to$ 51.4\% ($-5.7$\pp)
  \item Phase~15: Gated CARE-ALERT $\to$ 57.1\% (baseline restored)
\end{itemize}

% =====================================================================
\section{Discussion}\label{sec:discussion}
% =====================================================================

\subsection{The CRDT Unification}

The most significant theoretical contribution is the unification of
dialogue-level and memory-level synchronization under the same
mathematical framework.  \care Resolve operates at two temporal scales:
\textbf{intra-session} (MEANING-CRDT: merging agent positions in
real-time) and \textbf{inter-session} (MEMORY-CRDT: merging memory
states across time).

This suggests a general principle: weighted convergence under conflict
is a universal primitive for multi-agent coordination, applicable from
millisecond-level token generation to month-level knowledge
accumulation.

\subsection{Cognitive Bias in LLM Ensembles}

The hypercorrection discovery has implications beyond memory systems.
If LLM ensembles exhibit human-like cognitive biases under pre-session
priming, other known human biases may also manifest: anchoring effects
(confirmed in Phases~20--21), groupthink, and confirmation bias.

\subsection{Design Principles for AI Memory}

\begin{enumerate}
  \item \textbf{Create liberally, inject surgically}: Auto-log
    everything; inject only on detected divergence.
  \item \textbf{Operational over moralistic}: Tell agents what IS
    happening, not what they SHOULD do.
  \item \textbf{Current state only}: Never transmit historical
    failures or trust scores.
  \item \textbf{Gated intervention}: Monitor $\to$ Detect $\to$
    Intervene beats blanket pre-loading.
\end{enumerate}

% =====================================================================
\section{Limitations}\label{sec:limitations}
% =====================================================================

\begin{enumerate}
  \item $N = 40$ (35 unseen) --- insufficient power for small effects.
  \item Single agent trio tested --- only Alpha configuration.
  \item Non-significant primary accuracy result.
  \item Keyword retrieval (MVP) --- no semantic similarity.
  \item No consolidation implemented (Episodic $\to$ Semantic).
  \item MD thresholds (0.5, 0.8) set \textit{a priori} without
    optimization.
  \item Single run --- no repeated measurements.
  \item Cooperative-only scope --- no adversarial memory poisoning.
\end{enumerate}

% =====================================================================
\section{Conclusion}\label{sec:conclusion}
% =====================================================================

\mnemo extends the \arrival Protocol with persistent memory.  The
four-layer architecture provides structured knowledge organization
while MEMORY-CRDT~v1.1 ensures consistent synchronization using the
same mathematics that powers real-time dialogue convergence.

The three-phase experimental progression reveals fundamental principles:
\begin{itemize}
  \item \textbf{Phase~13} (Baseline): 57.1\% on GPQA Diamond.
  \item \textbf{Phase~14} (Global injection): $-5.7$\pp degradation
    $\to$ hypercorrection discovered.
  \item \textbf{Phase~15} (Gated CARE-ALERT): Baseline restored;
    first significant \care improvement ($p = 0.042$); 17.5\% alert
    rate.
\end{itemize}

Total project cost across all phases: \$1.17.  The entire investigation
--- from architecture design through negative result discovery to
validated real-time intervention --- cost less than \$1.20.

% =====================================================================
% Bibliography
% =====================================================================

\begin{thebibliography}{9}

\bibitem{kelevra2026arrival}
Kelevra, M. (2025--2026).
\newblock ARRIVAL Protocol: Cross-Architecture AI-to-AI Coordination
  Through Structured Semantic Atoms.
\newblock \textit{Preprint.}
\newblock DOI: \href{https://doi.org/10.5281/zenodo.18893515}%
  {10.5281/zenodo.18893515}.

\bibitem{kelevra2026crdt}
Kelevra, M. (2026).
\newblock MEANING-CRDT v1.1: Mathematical Foundation for Semantic
  Conflict Resolution.
\newblock DOI: \href{https://doi.org/10.5281/zenodo.18702383}%
  {10.5281/zenodo.18702383}.

\bibitem{du2024debate}
Du, Y., Li, S., Torralba, A., Tenenbaum, J.~B., \& Mordatch, I. (2024).
\newblock Improving Factuality and Reasoning in Language Models through
  Multiagent Debate.
\newblock \textit{ICML 2024.}

\bibitem{lewis2020rag}
Lewis, P., Perez, E., Piktus, A., et al. (2020).
\newblock Retrieval-Augmented Generation for Knowledge-Intensive NLP
  Tasks.
\newblock \textit{NeurIPS 2020.}

\bibitem{packer2023memgpt}
Packer, C., Wooders, S., Lin, K., et al. (2023).
\newblock MemGPT: Towards LLMs as Operating Systems.
\newblock \textit{arXiv:2310.08560.}

\bibitem{park2023agents}
Park, J.~S., O'Brien, J.~C., Cai, C.~J., et al. (2023).
\newblock Generative Agents: Interactive Simulacra of Human Behavior.
\newblock \textit{UIST 2023.}

\bibitem{rein2024gpqa}
Rein, D., Hou, B.~L., Stickland, A.~C., et al. (2024).
\newblock GPQA: A Graduate-Level Google-Proof Q\&A Benchmark.
\newblock \textit{arXiv:2311.12022.}

\bibitem{shapiro2011crdt}
Shapiro, M., Preguica, N., Baquero, C., \& Zawirski, M. (2011).
\newblock Conflict-free Replicated Data Types.
\newblock \textit{SSS 2011.}

\bibitem{tulving1972memory}
Tulving, E. (1972).
\newblock Episodic and Semantic Memory.
\newblock In: \textit{Organization of Memory.} Academic Press.

\end{thebibliography}

% =====================================================================
\appendix
\section{Memory Store Schema}

\begin{lstlisting}[caption={Memory dataclass definitions.}]
@dataclass
class EpisodicMemory:
    id: str              # Content-addressed hash
    session_id: str      # Source session identifier
    timestamp: str       # ISO 8601
    task: str            # Task description
    models: List[str]    # Participating models
    outcome: Dict        # {accuracy, care_resolve}
    care_resolve: float  # CARE Resolve score (0.0-1.0)
    meaning_debt: float  # Meaning Debt score
    key_insight: str     # Human-readable summary
    atoms_used: List     # @-atoms observed
    ttl_days: int = 30

@dataclass
class ProceduralMemory:
    id: str
    strategy_name: str
    task_type: str
    description: str
    success_rate: float
    n_trials: int

@dataclass
class SemanticMemory:
    id: str
    domain: str
    content: str
    confidence: float
    evidence_count: int

@dataclass
class MetaMemory:
    id: str
    agent_id: str
    trust_score: float
    domain_scores: Dict
    total_sessions: int
\end{lstlisting}

\section{CARE-ALERT Format Example}

\begin{lstlisting}[language={},caption={Example CARE-ALERT injection.}]
@CARE.ALERT [Level 1 -- Divergence Detected]
After Round 2, agents hold different positions:
- Alpha (GPT-4o): selected answer C
- Beta (DeepSeek-V3): selected answer A
Interim Meaning Debt: 0.72

Please address divergence with @EVIDENCE and
@COUNTER. Consider whether Beta's position has
merit before finalizing.
\end{lstlisting}

\section{MEMORY-CRDT Formal Properties}

\textbf{Commutativity}: For any memories $m_1, m_2$ with the same ID:
$\max(\mu(m_1), \mu(m_2)) = \max(\mu(m_2), \mu(m_1))$. \hfill$\square$

\textbf{Associativity}: For any three stores $A, B, C$:
$\text{merge}(\text{merge}(A,B), C) = \text{merge}(A, \text{merge}(B,C))$
because $\max$ is associative and set union is associative. \hfill$\square$

\textbf{Idempotency}: For any store $A$:
$\text{merge}(A, A) = A$ because for every memory $m \in A$,
$\max(\mu(m), \mu(m)) = \mu(m)$, and $A \cup A = A$. \hfill$\square$

\end{document}
